\addchap*{Abstract}
\markboth{Abstract}{Abstract}

About one gym member in three leaves each year, and most stop attending weeks before they cancel. Retention tools flag members at risk but leave the decision that matters (whom to contact this week, with which offer, within a fixed budget and limited staff hours) to guesswork. AnchorAI makes that decision and explains it. A calibrated ensemble predicts churn on real members; a forward-chaining rule base encodes the gym's playbooks and guardrails; case-based reasoning estimates what each offer would achieve for each member; a genetic algorithm, benchmarked against an exact CP-SAT solver, builds the weekly plan; a knowledge graph adds signals the churn model cannot see, such as a departed instructor; and a grounded assistant answers staff questions from the system's own evidence. Every recommendation carries its reasons, and managers approve, override and audit the plan in one web application.

Segmentation and churn prediction are evaluated on 4,000 real members of a public dataset. Decisions are evaluated on a documented synthetic scenario layer whose true offer effects are known, so plans can be graded against a truth no planner sees; their gains show that the method works when effects are of this kind, not how large they are in a real gym.

\vspace{1.2\baselineskip}
\noindent
\keyfigure{0.908}{churn F1 on 800 unseen real members (PR-AUC 0.973)}\hfill
\keyfigure{+22--138\%}{more members kept than contacting by risk or by schedule, at equal budget}\hfill
\keyfigure{0.24\%}{largest gap of the genetic algorithm to the proven optimum}\hfill
\keyfigure{96.2\%}{of the assistant's judged sentences supported by its evidence}

\vspace{1.4\baselineskip}
\noindent{\small\sffamily\textcolor{navy}{\textbf{Keywords}}\enskip churn prediction · uplift modelling · case-based reasoning · genetic algorithms · constraint programming · rule-based systems · knowledge graphs · retrieval-grounded assistants}
\clearpage
